% SOURCE AUTHORITY:
% expose_VIII_body.tex lines 1607--1696; scan indices 280--284;
% source folios 269--273.
% Direct scan comparison: physical PDF pages 281--284 in overlapping
% 1100-dpi bands.  The prose and displayed formulas are fully legible at
% that resolution; no 5000--9000-dpi ambiguity escalation was needed.
% Source disposition in the proof of Corollary 4.8: the print displays
% Biext^i(P,M;G_m).  This is retained.  Together with (4.7.1), applied with
% Q=M, those vanishings give the Biext^i(A,M;G_m) vanishings required by
% the exact sequence for 0 -> M -> Q -> B -> 0.
% Hard stop before 4.11, source line 1701 / scan index 284.

\medskip
\noindent\underline{Proposition 4.7.}  \underline{Let $P$ and $Q$
be two smooth connected algebraic groups over the field $k$.
Suppose that $P$ is an extension of an abelian scheme $A$ by a smooth
connected affine algebraic group $L$ admitting a composition series
each of whose factors is isomorphic to a torus or to $\underline
G_a$ (a condition always satisfied when $k$ is perfect).  Then the
functor considered in 1.3 is an equivalence of categories; that is,
every birigidified line bundle $\mathbb L$ on $P\times Q$ comes from
a biextension of $(P,Q)$.  Moreover, the functor between the rigid
categories}
\[
 \operatorname{BIEXT}(A,Q;\underline G_m)
 \longrightarrow
 \operatorname{BIEXT}(P,Q;\underline G_m)
\]
\underline{is an equivalence of categories; in other words, it
induces an isomorphism of groups}
\begin{equation}\tag{4.7.1}
 \Biext^1(A,Q;\underline G_m)
 \xrightarrow{\ \sim\ }
 \Biext^1(P,Q;\underline G_m).
\end{equation}

For the first assertion, descent and 4.3 again reduce us to the case
where $k$ is separably closed.  Proceeding as in the proof of 4.6,
we deduce from 4.5 and the hypothesis on $L$ that $\mathbb L$ is
isomorphic to the pullback of a birigidified line bundle $\mathbb M$
on $A\times Q$.  For the first assertion of 4.7 we are thus reduced
to the case $A=P$.  The line bundle $\mathbb L$ then defines a
morphism of pointed schemes
\[
 f:Q\longrightarrow\underline{\Pic}_{P/k},
\]
which necessarily factors through the dual abelian scheme $P'$ of
$P$.  Let $\mathbb L'$ be the pullback of the Weil sheaf on
$P\times P'$ along
\[
 \id_P\times f:P\times Q\longrightarrow P\times P'.
\]
This is a birigidified line bundle which, by the definition of $f$,
differs from $\mathbb L$ only by the pullback of a line bundle
$\mathbb M$ on $P$.  Since both $\mathbb L'$ and $\mathbb L$ have
trivial restriction to $P\times e_Q$, it follows that $\mathbb M$ is
trivial and hence that $\mathbb L\simeq\mathbb L'$.  We noted in VII
2.9.4 that divisorial correspondences on a product of two abelian
schemes arise from biextensions.  Consequently $\mathbb L'$, and
therefore $\mathbb L$, comes from a biextension.

Returning to the general case, where $P$ is an extension of $A$ by
$L$, it remains to prove that (4.7.1) is an isomorphism.  This follows
from the exact sequence of the $\Biext^i$, for $i=0,1$, associated
with
\[
 0\longrightarrow L\longrightarrow P\longrightarrow A
 \longrightarrow0,
\]
and from the relations
\[
 \Biext^0(L,Q;\underline G_m)=0,
 \qquad
 \Biext^1(L,Q;\underline G_m)=0.
\]
The first is a special case of the rigidity assertion in 4.3, and
the second is precisely the first case of 4.6.

\medskip
\noindent\underline{Corollary 4.8.}  \underline{Under the hypotheses
of 4.7, suppose that $Q$ is an extension of an algebraic group $B$
(necessarily smooth and connected) by a smooth connected affine
algebraic group $M$.  Then the functor between the rigid categories}
\[
 \operatorname{BIEXT}(A,B;\underline G_m)
 \longrightarrow
 \operatorname{BIEXT}(P,Q;\underline G_m)
\]
\underline{is an equivalence of categories; in other words, it
induces an isomorphism of groups}
\begin{equation}\tag{4.8.1}
 \Biext^1(A,B;\underline G_m)
 \xrightarrow{\ \sim\ }
 \Biext^1(P,Q;\underline G_m).
\end{equation}

In view of (4.7.1), it remains to prove that
\[
 \Biext^1(A,B;\underline G_m)
 \xrightarrow{\ \sim\ }
 \Biext^1(A,Q;\underline G_m)
\]
is an isomorphism.  As in the proof of (4.7.1), this follows from the
exact sequence of the $\Biext^i$ associated with
\[
 0\longrightarrow M\longrightarrow Q\longrightarrow B
 \longrightarrow0,
\]
together with the relations
\[
 \Biext^0(P,M;\underline G_m)=0,
 \qquad
 \Biext^1(P,M;\underline G_m)=0.
\]
Indeed, (4.7.1), applied with $M$ in place of $Q$, identifies these
groups with the corresponding groups having $A$ in the first
variable, as required by the preceding exact sequence.  The second
vanishing displayed above is a special case of the second case of
4.6, after interchanging the two factors.

\medskip
\noindent\underline{Remark 4.9.}  Let $P$ and $Q$ be two smooth
connected algebraic groups over a \underline{perfect} field $k$.
It is then known [3] that each is an extension of an abelian scheme
by a smooth connected affine group, the latter being the product of
a torus with an iterated extension of copies of $\underline G_a$
(cf. the reference given in 4.6):
\[
 0\longrightarrow L\longrightarrow P\longrightarrow A
 \longrightarrow0,
 \qquad
 0\longrightarrow M\longrightarrow Q\longrightarrow B
 \longrightarrow0.
\]
Thus 4.8 applies and shows that the theory of biextensions of $(P,Q)$
by $\underline G_m$ reduces to the analogous (and much more
classical!) theory for the pair $(A,B)$ of abelian varieties.

\medskip
\noindent\underline{Corollary 4.10.}  \underline{Let $P$ and $Q$ be
two smooth connected algebraic groups over $k$, let $L$ be a smooth
connected affine algebraic subgroup of $P$, and let $E$ be a
biextension of $(P,Q)$ by $\underline G_m$.  By 2.2, $E$ induces a
pairing}
\[
 T_\ell(P)\times T_\ell(Q)\longrightarrow T_\ell(\underline G_m).
\]
\underline{Then the induced pairing}
\[
 T_\ell(L)\times T_\ell(Q)\longrightarrow T_\ell(\underline G_m)
\]
\underline{is zero.}

Indeed, we may plainly suppose that $k$ is algebraically closed.  We
are then under the hypotheses of 4.6, which imply that the
restriction of $E$ to $(L,Q)$ is the trivial biextension.  The
asserted conclusion follows, since the induced pairing in question
is the pairing defined by this restriction.
